\chapter{Nouvelles et événements}\label{ch:s1:news-and-events}

Une dépêche arrive sur le fil à 10:32:07. En quelques millisecondes, les programmes qui analysent le flux ont déjà réagi, et l'essentiel du mouvement du cours est joué
avant qu'une personne ait fini de lire la première phrase. Il reste la part que le marché met des jours à admettre~: Tetlock, Saar-Tsechansky et Macskassy ont montré
que les cours réagissent insuffisamment, brièvement, aux mots négatifs des nouvelles propres à une entreprise, et Hirshleifer, Lim et Teoh ont montré que la réaction à une surprise de résultats
est plus faible, et la dérive qui la suit plus forte, les jours où beaucoup d'autres entreprises publient les leurs. Sur le flux de nouvelles synthétique de ce chapitre, un trader qui agit en une milliseconde
dispose encore de 96.7~\% du mouvement~; à un dixième de seconde, de 70.1~\%~; après une seconde, il ne reste que les 28.7~\% que les machines laissent aux
jours suivants. Un trader quotidien qui lit le ton de chaque élément à l'aide d'un modèle de texte imparfait obtient un ratio de Sharpe de 3.2 sur ce reliquat, davantage les jours d'actualité chargée que les jours calmes.
La construction s'appelle \texttt{firm.newsevent}.

\section{Dépêches lisibles par machine}

\begin{definition}[Dépêches lisibles par machine]\label{def:s1:news-and-events:mrn}
Les \emph{dépêches lisibles par machine}\index{dépêches lisibles par machine} (machine-readable news) sont un flux d'informations diffusé sous une forme structurée, que des programmes peuvent exploiter sans qu'une personne ait à le lire~:
chaque élément porte un horodatage, les titres concernés, une catégorie et, souvent, des scores de pertinence, de nouveauté et de ton, calculés par le fournisseur ou
par le modèle de texte propre à l'utilisateur.
\end{definition}

Une stratégie fondée sur les nouvelles transforme les éléments en positions en trois étapes. Elle associe chaque élément aux titres qu'il concerne et écarte ceux qui ne les concernent pas.
Elle évalue la direction et l'intensité de chaque élément, à partir d'un dictionnaire de mots positifs et négatifs, dans sa version la plus simple, ou d'un modèle
de langage entraîné sur les réactions passées. Enfin, elle décide du moment d'agir~: dans la première seconde, lorsque le mouvement reste à venir mais que seules les machines s'en disputent, ou plus tard, sur la part
que le marché met plus de temps à absorber. La littérature publiée porte surtout sur le second cas. Tetlock a mesuré le pessimisme d'une chronique quotidienne du \emph{Wall Street
Journal} et a constaté qu'un pessimisme élevé annonce une pression baissière sur les cours, suivie d'un retour vers les fondamentaux~: le ton y apparaît comme un sentiment de traders bruitistes.
Tetlock, Saar-Tsechansky et Macskassy ont appliqué cette mesure aux articles propres à une entreprise et ont constaté que la proportion de mots négatifs annonce de faibles résultats,
et que les cours y réagissent insuffisamment, brièvement, surtout pour les articles portant sur les fondamentaux~: le ton y apparaît comme une information que le marché met du temps à absorber.

La série du chapitre (\texttt{firm.newsevent}) couvre huit ans d'éléments concernant 1\,000 actions, soit 40.8 par jour en moyenne (82\,351 au total), avec une intensité de l'information
qui varie d'un jour à l'autre. Chaque élément possède un ton, c'est-à-dire le mouvement total qu'il entraînera, avec un écart type de 2~\%~; un niveau d'attention, qui diminue
lorsque d'autres éléments arrivent le même jour ($1/(1+0.03(k-1))$ pour $k$ éléments, avec une médiane de 0.42)~; et une heure d'arrivée, uniforme sur la séance. Une part du mouvement, qui va de la moitié
sans attention à la totalité, avec attention pleine, se produit en une fraction de seconde, avec une constante de temps de 0.2 seconde~; le reste se répartit
régulièrement sur les cinq jours suivants.

\section{Événements programmés}

Les annonces de résultats, les déclarations de dividendes, les changements d'indices et les publications économiques surviennent à des dates connues, et le chapitre~7 a étudié la dérive qui suit la première
d'entre elles au moyen d'une \omterm{def:s1:earnings:event}{étude d'événement}. Pour le trading sur nouvelles, le calendrier change deux choses. La réaction peut être préparée~: le consensus est connu, la surprise tient en un seul nombre,
et l'ordre peut être préparé avant la publication, ce qui place la première seconde entre les mains de ceux qui disposent des flux les plus rapides. Et la distraction est prévisible~:
Hirshleifer, Lim et Teoh ont mesuré la charge d'information de l'investisseur par le nombre d'annonces du jour, et ont constaté que la
réaction immédiate du cours et du volume à une surprise est beaucoup plus faible, et la dérive postérieure à l'annonce beaucoup plus forte, lorsque beaucoup d'autres entreprises publient le même jour
au point de procurer à une stratégie de trading des alphas substantiels. Le flux synthétique intègre leur mécanisme dans sa fonction d'attention.

\begin{definition}[Fenêtre de réaction aux nouvelles]\label{def:s1:news-and-events:window}
Une \emph{fenêtre de réaction aux nouvelles}\index{fenêtre de réaction aux nouvelles} est l'intervalle qui suit un élément d'information, pendant lequel une stratégie mesure ou négocie le mouvement du cours, depuis
l'horodatage de l'élément, augmenté de la latence de la stratégie, jusqu'à la fin de la période de détention~; la part du mouvement total de l'élément qui tombe à l'intérieur
de cette fenêtre est ce que la stratégie peut capturer.
\end{definition}

\section{Événements non programmés et suspensions de cotation}

Les nouvelles non programmées (fusions, avertissements sur résultats, mesures réglementaires, accidents) surviennent à tout moment, et les bourses disposent d'un outil pour les éléments les plus importants~:
elles suspendent la cotation jusqu'à la diffusion de l'information.

\begin{definition}[Suspension de cotation]\label{def:s1:news-and-events:halt}
Une \emph{suspension de cotation}\index{suspension de cotation} (trading halt) est l'arrêt de la négociation d'un titre, ordonné par la bourse de cotation ou par le régulateur, dans l'attente d'une nouvelle, après que la nouvelle a été diffusée,
ou après une variation de cours trop forte pour que le marché l'absorbe en continu~; la négociation reprend, en général par une enchère de réouverture, lorsque
la suspension est levée.
\end{definition}

\begin{dated}{2026-09}{Codes de suspension sur Nasdaq}\label{dat:s1:news-and-events:halts}
Nasdaq publie un code avec chaque suspension. Parmi eux figurent T1, «~Halt -- News Pending~»~; T2, «~Halt -- News Released~»~; T3, les heures de l'information et de la reprise~; T5, une
suspension de négociation d'un seul titre~; LUDP, une suspension de négociation pour volatilité sous le régime de Limit Up--Limit Down~; et MWC1 à MWC3, les trois niveaux du
coupe-circuit de marché. Le livre~1, chapitre~31 décrit les bandes de Limit Up--Limit Down et les coupe-circuits.
\end{dated}

Une suspension modifie qui capte le mouvement. Lorsqu'un titre est suspendu, personne ne peut le négocier en continu, de sorte que le mouvement se produit à l'enchère de réouverture, où
les ordres de chacun se rencontrent à un seul prix~: la vitesse dans la première seconde ne rapporte rien, et ce qui compte est la manière dont les ordres sont placés dans l'enchère et ce qui se passe
après celle-ci. Le flux synthétique suspend chaque élément qui fait varier son titre de plus de 5~\% (1.18~\% des éléments) pendant cinq minutes, et place la part immédiate du mouvement
à la réouverture. Un trader rapide ne capte rien du mouvement immédiat d'un élément suspendu~; la dérive qui suit est la même que pour n'importe quel autre élément.
Savoir si les prix de réouverture surréagissent ou sous-réagissent est une question empirique que le modèle du chapitre ne tranche pas~; la fiche de la
stratégie post-suspension indique ce qu'il faut tester.

\section{Vitesse de réaction et capture du mouvement}

Le \Cref{lst:s1:news-and-events:remaining} présente l'ensemble du modèle de réaction~: la part du mouvement d'un élément qui reste à capter pour un trader qui agit à une latence donnée.
En pondérant les éléments par l'ampleur de leur mouvement, on obtient la part du mouvement total du flux qui reste disponible.

\begin{center}
\small
\begin{tabular}{@{}lcccccc@{}}
\toprule
latence & 1 ms & 10 ms & 100 ms & 1 s & 1 heure & clôture suivante\\
\midrule
part du mouvement encore disponible & 96.7~\% & 93.7~\% & 70.1~\% & 29.2~\% & 28.7~\% & 28.7~\%\\
P\&L par élément après 10 pb & 144 & 139 & 102 & 36 & 36 & 36\\
\bottomrule
\end{tabular}
\end{center}

La courbe (\cref{fig:s1:news-and-events:latency}) présente deux régimes. En deçà d'une seconde, la part diminue rapidement, car les machines prennent la part immédiate
en quelques constantes de temps~: un trader en retard d'un dixième de seconde a perdu près d'un tiers du mouvement au profit des plus rapides. Au-delà d'une seconde, elle est plate~: les machines ont terminé, et
ce qui reste, 28.7~\%, est la part que l'attention du marché laisse aux jours suivants. Les P\&L supposent que le trader connaît la direction de
chaque élément, comme le font en pratique les traders les plus rapides pour les éléments les plus simples~; ils mesurent la valeur de la fenêtre, et non celle d'une stratégie.

\begin{omfigure}
\begin{tikzpicture}
\begin{axis}[width=10cm,height=5.4cm,xmode=log,xlabel={\small latence après l'élément (secondes)},ylabel={\small part du mouvement encore disponible (\%)},
  tick label style={font=\footnotesize},xmin=0.001,xmax=100,ymin=0,ymax=100,grid=major,grid style={gray!25},legend pos=north east,
  legend style={font=\scriptsize},table/col sep=comma]
\addplot[omThm,very thick] table[x=latency,y=share_pct]{figdata/strategies-1/17-news-and-events/latency.csv};
\addplot[black!40,dashed] coordinates {(0.001,28.7) (100,28.7)};
\legend{flux de nouvelles synthétique,part laissée à la clôture suivante}
\end{axis}
\end{tikzpicture}

\omcaption{La part du mouvement total du flux de nouvelles synthétique qui reste à capter pour un trader agissant à une latence donnée après chaque élément, pondérée par l'ampleur
du mouvement de chaque élément. Les machines prennent la part immédiate en une fraction de seconde~; la ligne en pointillés figure la dérive restant à la clôture suivante. Données~:
\texttt{s1\_news.capture}.}\label{fig:s1:news-and-events:latency}
\end{omfigure}

Un trader quotidien ne peut pas rivaliser pendant la première seconde, mais il dispose du reste. Le trader quotidien du chapitre lit le ton de chaque élément à l'aide d'un modèle de texte dont l'erreur
vaut 1.5 fois l'écart type du ton~; il trouve ainsi la bonne direction dans 68.9~\% des cas, achète ou vend à la clôture suivante dans le sens de sa lecture et
conserve la position pendant cinq jours, avec un bruit quotidien de 2~\% sur chaque titre et dix points de base pour un aller-retour (\cref{lst:s1:news-and-events:daily}).

\begin{center}
\small
\begin{tabular}{@{}lc>{\centering\arraybackslash}p{4.4cm}c@{}}
\toprule
éléments & ensemble & jours calmes (attention au-dessus de la médiane) & jours chargés\\
\midrule
éléments négociés & 82\,351 & 39\,735 & 42\,616\\
dérive captée par élément (pb) & 25.4 & 20.8 & 29.5\\
rendement annuel après coûts & 7.4~\% & 6.3~\% & 8.3~\%\\
ratio de Sharpe après coûts & 3.24 & 1.80 & 2.74\\
\bottomrule
\end{tabular}
\end{center}

Les jours chargés laissent davantage de dérive par élément, comme chez Hirshleifer, Lim et Teoh~: l'attention du marché se répartit sur davantage de nouvelles, si bien qu'une moindre part de chaque
mouvement se produit d'un seul coup. Chacun des portefeuilles issus de la scission détient environ la moitié des éléments, de sorte que leurs ratios de Sharpe sont inférieurs à celui du portefeuille entier, qui diversifie sur deux fois plus d'éléments. Le
ratio de Sharpe de 3.2 est élevé parce que le modèle introduit proprement la sous-réaction et que les erreurs de lecture sont indépendantes de tout le reste~; les effets publiés
se mesurent en points de base par nouvelle, et les coûts liés à la négociation de nombreuses petites positions les absorbent rapidement.

\section{Fiches de stratégie}

\begin{strategyfile}{Sentiment des titres}\label{strat:s1:news-and-events:sentiment}
\sfield{Qui paie, et pourquoi} Les investisseurs qui intègrent lentement l'information contenue dans le texte des nouvelles, et les noise traders qui font bouger les cours sur le seul ton.
\sfield{Instruments et plateformes} Actions liquides dont l'actualité propre aux sociétés est fréquente~; contrats à terme sur indices pour le ton de l'ensemble du marché.
\sfield{Signal} Le ton des nouvelles propres aux sociétés, issu d'un dictionnaire ou d'un modèle de texte entraîné, filtré selon la pertinence et la nouveauté.
\sfield{Dimensionnement et exécution} Positions prises après les premières secondes, dans le sens du ton, conservées quelques jours~; nombreuses petites positions.
\sfield{Coûts} Flux de nouvelles et traitement de texte~; rotation à chaque nouvelle.
\sfield{Comment elle s'éteint} De meilleurs modèles de texte partout~; la dérive absorbée plus vite.
\sfield{Horizon, capacité, infrastructure} Quelques jours~; capacité limitée par le nombre de nouvelles et leur liquidité~; une archive de nouvelles horodatées.
\sfield{Backtest en toute honnêteté} Les nouvelles telles qu'elles sont parvenues, avec les horodatages propres au flux, et non la date de publication de l'article~; aucun score de ton révisé.
\sfield{Sources} Tetlock (2007)~; Tetlock, Saar-Tsechansky et Macskassy (2008)~: les cours sous-réagissent brièvement aux mots négatifs.
\end{strategyfile}

\begin{strategyfile}{Dérive sur événements programmés}\label{strat:s1:news-and-events:scheduled}
\sfield{Qui paie, et pourquoi} Les investisseurs qui sous-réagissent aux nouvelles programmées, d'autant plus que de nombreuses annonces se disputent leur attention.
\sfield{Instruments et plateformes} Actions dont des annonces figurent au calendrier.
\sfield{Signal} La surprise par rapport au consensus, lue dans la publication dès son arrivée.
\sfield{Dimensionnement et exécution} Positions prises à la clôture du jour d'annonce, dans le sens de la surprise, conservées de quelques jours à quelques semaines.
\sfield{Coûts} Modérés~; négociation proche des clôtures encombrées.
\sfield{Comment elle s'éteint} Des réactions plus rapides des machines laissent moins de dérive.
\sfield{Horizon, capacité, infrastructure} De quelques jours à quelques semaines~; un calendrier des événements et une base de consensus.
\sfield{Backtest en toute honnêteté} Heures de publication à la seconde près~; consensus tel qu'il était avant la publication.
\sfield{Sources} L'\omterm{def:s1:earnings:event}{étude d'événement} du chapitre~7~; Hirshleifer, Lim et Teoh (2009).
\end{strategyfile}

\begin{strategyfile}{Réouverture après suspension}\label{strat:s1:news-and-events:halt}
\sfield{Qui paie, et pourquoi} Les traders qui doivent négocier à la réouverture quel que soit le cours, et une enchère de réouverture dont le cours peut dépasser son niveau d'équilibre.
\sfield{Instruments et plateformes} Actions suspendues pour cause de nouvelles ou de volatilité~; l'enchère de réouverture.
\sfield{Signal} Les nouvelles publiées pendant la suspension, le cours indicatif et le déséquilibre de l'enchère, et le mouvement depuis la suspension.
\sfield{Dimensionnement et exécution} Ordres dans l'enchère de réouverture ou juste après, à contre-sens d'un dépassement ou dans le sens d'une dérive~; tailles réduites, car la liquidité est faible.
\sfield{Coûts} Fourchettes larges après la réouverture~; risque d'une nouvelle suspension.
\sfield{Comment elle s'éteint} La concurrence pour la liquidité à la réouverture.
\sfield{Horizon, capacité, infrastructure} De quelques minutes à quelques jours~; le flux des suspensions et les cours indicatifs de l'enchère.
\sfield{Backtest en toute honnêteté} Suspensions et heures de reprise tirées des registres de la bourse~; cours d'enchère, et non les premières transactions qui suivent.
\sfield{Sources} Aucun chiffre de performance vérifié~; les codes de suspension de l'encadré daté.
\end{strategyfile}

\begin{strategyfile}{Sous-réaction aux nouvelles peu remarquées}\label{strat:s1:news-and-events:attention}
\sfield{Qui paie, et pourquoi} Les investisseurs distraits, qui réagissent moins aux nouvelles les jours où d'autres informations se disputent davantage leur attention.
\sfield{Instruments et plateformes} Actions avec des nouvelles les jours chargés.
\sfield{Signal} Le sens de l'élément, pondéré par le volume de nouvelles du jour (le nombre d'annonces ou de nouvelles publiées le même jour).
\sfield{Dimensionnement et exécution} Positions plus importantes les jours chargés~; conservées quelques jours.
\sfield{Coûts} Comme pour le sentiment des titres.
\sfield{Comment elle s'éteint} Des machines qui ne se laissent pas distraire.
\sfield{Horizon, capacité, infrastructure} Quelques jours~; un décompte des nouvelles du jour.
\sfield{Backtest en toute honnêteté} Le volume de nouvelles tel qu'il était connu au moment de la transaction.
\sfield{Sources} Hirshleifer, Lim et Teoh (2009)~: dérive plus forte et alphas substantiels~; la répartition par attention de ce chapitre.
\end{strategyfile}

\section{Tutoriel~: avant qu'un humain ne le lise}\label{tut:s1:news-and-events:tut}

\begin{tutorial}
\textbf{Objectif.} Générer un flux de nouvelles avec un ton et une attention injectés, mesurer la part de son mouvement qui reste à venir à chaque latence, et négocier ce qui reste
à la clôture suivante. \textbf{État final~:} les deux tableaux et la \cref{fig:s1:news-and-events:latency}.
\begin{tutsteps}
\item \textbf{Le modèle de réaction}~: la part immédiate selon l'attention, et ce qui reste à une latence donnée~; les éléments suspendus n'ont rien à capter.
\omcode{code/firm/newsevent/firm_newsevent.py}{55}{68}{La part immédiate et le mouvement restant à venir à une latence donnée.}\label{lst:s1:news-and-events:remaining}
\item \textbf{Le trader quotidien}~: une lecture bruitée du ton de chaque élément, entrée à la clôture suivante, conservation de la position pendant cinq jours.
\omcode{code/strategies-1/17-news-and-events/python/s1_news.py}{48}{68}{Le portefeuille du trader quotidien selon l'attention.}\label{lst:s1:news-and-events:daily}
\item \textbf{Exécuter} \texttt{stream()}, \texttt{capture(latency)} pour chaque latence et \texttt{daily(split)} pour les trois découpages, ainsi que \texttt{fig\_news.py}.
\end{tutsteps}
\textbf{Pour aller plus loin.} Faire varier l'erreur du modèle de texte selon la catégorie de l'élément~; faire dépasser son niveau d'équilibre à l'enchère de réouverture et négocier à contre-sens~; donner
au trader quotidien un signal pondéré par le volume de nouvelles du jour.
\end{tutorial}

\section{Construction~: un flux de nouvelles}\label{bld:s1:news-and-events:newsevent}

\begin{build}
\bfield{Rôle} Un flux de nouvelles synthétique avec ton, attention et suspensions, et la part du mouvement de chaque élément qui reste à venir à une latence donnée.
\bfield{Interface} \texttt{NewsConfig(\dots)}, \texttt{simulate\_news(days, stocks, cfg, rng)}, \texttt{immediate\_share(attention, cfg)},
\texttt{remaining(items, latency, cfg)} (latence \texttt{None} pour la clôture suivante).
\bfield{Règles} L'attention ne dépend que du nombre d'éléments du jour~; les éléments suspendus rouvrent au nouveau cours~; le flux est déterministe pour une graine donnée.
\bfield{Tests d'acceptation} \texttt{code/firm/newsevent/tests/}~: le seuil de suspension et la fonction d'attention~; la part restante pour des latences nulle, faible et
élevée.
\bfield{Pour aller plus loin} Scores de pertinence et de nouveauté~; éléments qui concernent plusieurs actions~; enchères de réouverture avec un déséquilibre.
\end{build}

\begin{omsources}
\item P.~C.~Tetlock, ``Giving content to investor sentiment: the role of media in the stock market'', \emph{Journal of Finance} 62(3), 2007.
\item P.~C.~Tetlock, M.~Saar-Tsechansky et S.~Macskassy, ``More than words: quantifying language to measure firms' fundamentals'', \emph{Journal of
Finance} 63(3), 2008.
\item D.~Hirshleifer, S.~S.~Lim et S.~H.~Teoh, ``Driven to distraction: extraneous events and underreaction to earnings news'', \emph{Journal of
Finance} 64(5), 2009.
\end{omsources}

\section{Exercices}

\begin{exercise}[$\star$]\label{exo:s1:news-and-events:1}
La partie immédiate d'un mouvement décroît avec une constante de temps de 0.2 seconde. Quelle part de celle-ci reste à venir pour un trader à 100 millisecondes, et à une
seconde~?
\end{exercise}

\begin{exercise}[$\star$]\label{exo:s1:news-and-events:2}
Un jour où 40 items sont publiés, quelle est l'attention portée à chaque item, et quelle part de son mouvement se produit immédiatement~? Quelle est cette part à l'attention médiane de 0.42~?
\end{exercise}

\begin{exercise}[$\star$]\label{exo:s1:news-and-events:3}
Pourquoi un trader rapide ne capte-t-il rien du mouvement immédiat d'un item soumis à une \omterm{def:s1:news-and-events:halt}{suspension de cotation}~?
\end{exercise}

\begin{exercise}[$\star\star$]\label{exo:s1:news-and-events:4}
Un modèle de texte lit un ton distribué normalement, auquel s'ajoute une erreur normale indépendante égale à 1.5 fois l'écart type du ton. À quelle fréquence le
modèle trouve-t-il la bonne direction~? (Pour un couple normal bivarié de corrélation $\rho$, les signes concordent avec la probabilité $\tfrac12 + \arcsin(\rho)/\pi$.)
\end{exercise}

\begin{exercise}[$\star\star$]\label{exo:s1:news-and-events:5}
Pourquoi la courbe de latence est-elle plate entre une seconde et la clôture suivante~?
\end{exercise}

\begin{exercise}[$\star\star$]\label{exo:s1:news-and-events:6}
Pourquoi les portefeuilles des journées calmes et ceux des journées chargées affichent-ils tous deux un ratio de Sharpe inférieur à celui du portefeuille complet~?
\end{exercise}

\begin{exercise}[$\star\star\star$]\label{exo:s1:news-and-events:7}
\emph{Programmation.} Fixer \texttt{READ} à 0 dans \texttt{s1\_news.py}, soit une lecture parfaite, et relancer \texttt{daily()}. Que devient le ratio de Sharpe, et que
mesure la différence~?
\end{exercise}

\begin{exercise}[$\star\star\star$]\label{exo:s1:news-and-events:8}
\emph{Trouver la faille.} «~Notre backtest échange chaque dépêche à la minute de son heure de publication sur le site et gagne 40 points de base par dépêche.~»
\end{exercise}

\section{Problème~: avant qu'une personne ne le lise}

\begin{problem}[{Problème du week-end --- la vie d'un titre}]\label{pb:s1:news-and-events:1}
Le flux synthétique de nouvelles du chapitre et le registre public.

\medskip\noindent\textbf{Partie I --- Les nouvelles comme données.}
\begin{enumerate}
\item Définir les \omterm{def:s1:news-and-events:mrn}{dépêches lisibles par machine} et la \omterm{def:s1:news-and-events:window}{fenêtre de réaction aux nouvelles}.
\item Énumérer les trois étapes qui mènent d'un item à une position.
\item Que Tetlock (2007) a-t-il établi au sujet du pessimisme médiatique~?
\item Que Tetlock, Saar-Tsechansky et Macskassy ont-ils établi au sujet des mots négatifs~?
\end{enumerate}

\medskip\noindent\textbf{Partie II --- Événements et suspensions.}
\begin{enumerate}[resume]
\item Qu'est-ce qui change lorsqu'un événement est programmé~?
\item Que Hirshleifer, Lim et Teoh ont-ils établi au sujet de la distraction~?
\item Définir une \omterm{def:s1:news-and-events:halt}{suspension de cotation} et nommer trois codes de suspension du Nasdaq.
\item Comment une suspension modifie-t-elle la personne qui capte le mouvement~?
\end{enumerate}

\medskip\noindent\textbf{Partie III --- Le flux.}
\begin{enumerate}[resume]
\item Décrire le flux synthétique~: items, ton, attention, réaction et suspensions.
\item Combien y a-t-il d'items, combien par jour, et quelle part est suspendue~?
\item Donner la part du mouvement encore à venir à 1 milliseconde, 100 millisecondes, une seconde et la clôture suivante.
\item Expliquer les deux régimes de la courbe de latence.
\end{enumerate}

\medskip\noindent\textbf{Partie IV --- Le verdict.}
\begin{enumerate}[resume]
\item Énoncer le \emph{résultat nommé}~: la part du mouvement lié aux nouvelles captée à chaque latence, et la dérive qui reste au trader quotidien.
\item Pourquoi les journées chargées laissent-elles davantage de dérive~?
\item Que coûte au trader quotidien l'erreur du modèle de texte~?
\item Pourquoi le ratio de Sharpe de 3.2 du trader quotidien constitue-t-il une borne supérieure pour les vraies nouvelles~?
\item Comment backtester honnêtement une stratégie fondée sur les nouvelles~?
\item Quelle fiche de stratégie exige l'infrastructure la plus rapide~?
\item Comment ce chapitre se rattache-t-il au chapitre~7~?
\item En une phrase~: que rapporte la vitesse dans le trading de nouvelles~?
\end{enumerate}
\end{problem}

\section{Questions d'entretien}

\begin{interviewq}[$\star$ \iqroles{chercheur}]\label{iq:s1:news-and-events:1}
Comment transformeriez-vous un flux de nouvelles en signal de trading~?
\end{interviewq}

\begin{interviewq}[$\star\star$ \iqroles{chercheur}]\label{iq:s1:news-and-events:2}
Votre backtest sur les nouvelles affiche des rendements élevés. Quels horodatages vérifieriez-vous en premier~?
\end{interviewq}

\begin{interviewq}[$\star\star$ \iqroles{trader}]\label{iq:s1:news-and-events:3}
Une action que vous détenez fait l'objet d'une \omterm{def:s1:news-and-events:halt}{suspension de cotation} alors qu'une nouvelle est en attente. Que faites-vous~?
\end{interviewq}

\begin{interviewq}[$\star\star$ \iqroles{développeur}]\label{iq:s1:news-and-events:4}
Concevez le chemin qui mène du flux d'un fournisseur de nouvelles à un ordre, pour une stratégie qui doit agir en dix millisecondes.
\end{interviewq}

\begin{interviewq}[$\star\star$ \iqroles{chercheur, trader}]\label{iq:s1:news-and-events:5}
Pourquoi un signal fondé sur les nouvelles pourrait-il mieux fonctionner les jours où les annonces de résultats sont nombreuses~?
\end{interviewq}

\begin{interviewq}[$\star\star\star$ \iqroles{chercheur}]\label{iq:s1:news-and-events:6}
L'impact immédiat d'un item décroît en $e^{-t/\tau}$ et une part $a$ du mouvement total est immédiate. Établir la part du mouvement encore à venir pour un trader à la
latence $L$, et la latence à laquelle la moitié de la part immédiate a disparu.
\end{interviewq}
